% Lịch sử: bản bàn giao tại 1b37300, trước biên tập P1 ngày 06/10/2026.
\chapter{Cài đặt PODEM}

\label{sec:p5-cai-dat-podem}

\section{Kiến trúc cài đặt}

Phần P5 triển khai lõi thuật toán PODEM trên mạch logic tổ hợp với mô hình
logic năm giá trị. Giao diện chính nhận một mạch \texttt{Circuit} và một
\texttt{Fault}, đồng thời hỗ trợ danh sách fault tương ứng với các khung
thời gian sau khi mạch tuần tự được unroll.

Luồng xử lý chính gồm các bước Imply, kiểm tra phát hiện fault, Objective,
Backtrace, gán Primary Input và Backtrack khi nhánh hiện tại thất bại.

\section{Logic năm giá trị}

Hệ thống sử dụng năm giá trị $0$, $1$, $X$, $D$ và $D'$.

Trong đó $D$ biểu diễn cặp giá trị good/faulty là $(1,0)$ và $D'$ biểu diễn
$(0,1)$. Giá trị $X$ biểu diễn trạng thái chưa xác định.

Pattern của Primary Input chỉ sử dụng các giá trị $0$, $1$ và $X$.
Giá trị $D$ và $D'$ được sử dụng trong quá trình implication để biểu diễn
sự khác nhau giữa mạch tốt và mạch lỗi.

Sau mỗi lần thay đổi pattern, mạch được implication lại theo thứ tự topo.

\section{Objective và Backtrace}

Khi fault chưa được kích hoạt, Objective đặt net lỗi về giá trị đối lập với
stuck-at value.

Sau khi fault được kích hoạt, Objective được chọn trên D-frontier nhằm lan
truyền giá trị $D$ hoặc $D'$ đến Primary Output.

Đối với các cổng NAND, NOR và NOT, giá trị Objective được đảo.
Đối với AND, OR và BUFF, giá trị Objective được giữ nguyên.

Đối với cổng XOR và XNOR, Objective được xử lý theo tính chẵn lẻ của các
đầu vào. Với XOR tương ứng $q=0$ và XNOR tương ứng $q=1$. Khi có một đầu vào
$X$, giá trị cần gán cho đầu vào đó được xác định từ giá trị mục tiêu và
giá trị trên rail tốt của các đầu vào đã biết. Khi có nhiều đầu vào $X$,
đầu vào $X$ đầu tiên theo thứ tự \texttt{Gate.inputs} được chọn để tiếp tục
Backtrace.

Backtrace đưa Objective từ một net trung gian về Primary Input chưa được
quyết định. Thứ tự lựa chọn được giữ ổn định theo thứ tự ngõ vào của gate.

\section{Backtrack}

Các quyết định Primary Input được lưu trong decision stack.

Khi nhánh hiện tại thất bại, quyết định được đảo sang giá trị còn lại và
thuật toán tiếp tục tìm kiếm. Biến đếm backtrack chỉ tăng khi thuật toán
thực sự chuyển sang một nhánh mới.

Sau khi backtrack, pattern được tạo lại từ decision stack và mạch được
implication lại trước khi tiếp tục. Trạng thái net và D-frontier được sử
dụng trong trace là trạng thái sau lần implication mới.

\section{D-frontier và điều kiện dừng}

D-frontier gồm các gate có đầu ra đang là $X$ và có đầu vào mang giá trị $D$
hoặc $D'$.

Thuật toán kết thúc thành công khi một Primary Output có giá trị $D$ hoặc
$D'$.

\section{Fault injection}

Stem fault tác động đến các fanout của net, trong khi branch fault chỉ tác
động đến branch được chỉ định.

Khi cấy stuck-at vào một giá trị logic năm giá trị, rail của mạch tốt được
giữ nguyên và rail của mạch lỗi được ép về stuck-at value.

Do đó, với các trạng thái đã chứa fault effect:

\[
D + SA0 \rightarrow D
\]

\[
D + SA1 \rightarrow 1
\]

\[
D' + SA0 \rightarrow 0
\]

\[
D' + SA1 \rightarrow D'
\]

Cách xử lý này tránh tạo ra fault effect $D/D'$ giả khi fault được cấy lặp
lại trên các vị trí hoặc các frame khác nhau.

\section{Hỗ trợ fault-frame}

Để tương thích với phần P4, PODEM có thể nhận danh sách fault tương ứng với
các frame của mạch sau khi unroll.

Khi sử dụng danh sách fault, \texttt{imply()} vẫn mô phỏng toàn bộ
fault-instance, trong khi từng fault-frame được lần lượt chọn làm target
cho Objective và Backtrace.

\section{Kết quả kiểm thử}

Với mạch c17 và fault 11/SA0, thuật toán tạo pattern

\[
X10XX
\]

với $0$ backtrack và trạng thái \texttt{DETECTED}.

Với ví dụ backtrack, thuật toán tạo pattern

\[
01
\]

với $1$ backtrack.

Các kiểm thử bổ sung cho XOR/XNOR bao gồm các tổ hợp logic năm giá trị,
nhiều đầu vào, lan truyền và triệt tiêu $D/D'$, cùng các trường hợp
Backtrace và PODEM theo đặc tả P3-v1.1.

Bộ kiểm thử hiện tại đạt:

\[
\boxed{262\ \text{passed}}
\]

Trong đó:

\[
\boxed{220\ \text{passed}}
\]

cho \texttt{tests/test\_logic.py}, và

\[
\boxed{32\ \text{passed}}
\]

cho \texttt{tests/test\_podem.py}.

\section{Tóm tắt}

Phần P5 đã triển khai các thành phần chính của PODEM gồm logic năm giá trị,
Objective, Backtrace, implication, D-frontier, Backtrack, fault injection
và trace.

Phần cài đặt cũng đã được bổ sung xử lý cho các cổng XOR/XNOR, fault effect
$D/D'$, fault-frame và trạng thái trace sau backtrack. Toàn bộ bộ kiểm thử
hiện tại của project đạt 262 test thành công.